% !TEX root = ../main.tex

\subsection{Algorithm choice}

The two models are LightGBM LambdaRank~\cite{ke2017lightgbm} and XGBoost \texttt{rank:ndcg}~\cite{chen2016xgboost}. Both are gradient-boosted decision tree implementations of pairwise / listwise ranking objectives, but they differ in two substantive ways: LightGBM uses histogram-based split finding with native categorical support and a leaf-wise tree growth strategy, while XGBoost uses level-wise growth with a regularised objective and approximate split finding. The two implementations therefore explore the feature space along noticeably different trajectories. Liu et al.~\cite{liu2013expedia} made the same diversity argument when justifying their thirty-base-learner ensemble: ensemble lift comes from disagreement, and two implementations of roughly the same ranking idea disagree enough to be worth blending. The rubric of the in-class assignment also asks for two distinct techniques to be compared; LightGBM and XGBoost meet this requirement while keeping the implementation effort manageable.

\subsection{Hyperparameter optimisation with Optuna TPE}

We use Optuna's TPE sampler~\cite{bergstra2011tpe,akiba2019optuna} on each model independently, with 50 trials per model and val NDCG@5 as the objective. The TPE sampler fits a probabilistic surrogate over the search space after a 10-trial random warmup and then concentrates subsequent trials in regions of high estimated NDCG. The choice over a single-axis manual grid is twofold: TPE finds non-obvious parameter interactions that a univariate sweep misses, and it amortises tuning effort across both models with one consistent protocol. Each trial fits up to 1000 trees with early stopping at 30 rounds on val NDCG@5. The seed is fixed at 42, the studies are deterministic, and the holdout is untouched.

\begin{table}[t]
\centering
\caption{Hyperparameter search spaces and Optuna-selected values for LightGBM LambdaRank and XGBoost \texttt{rank:ndcg}. Best trial val NDCG@5 in the LightGBM study reaches 0.40999 at 423 trees; XGBoost reaches 0.41223 in the refit (Section~\ref{subsec:ensemble}) at 694 boosting rounds.}
\label{tab:hyper}
\begin{tabular}{lll}
\toprule
\multicolumn{3}{l}{\textbf{LightGBM LambdaRank}} \\
\midrule
Hyperparameter & Search space & Selected \\
\midrule
learning\_rate & loguniform [0.01, 0.10] & 0.0565 \\
num\_leaves & int [15, 127] & 110 \\
max\_depth & int [4, 12] & 7 \\
min\_child\_samples & int [10, 200] & 74 \\
bagging\_fraction & uniform [0.5, 1.0] & 0.817 \\
feature\_fraction & uniform [0.5, 1.0] & 0.563 \\
lambda\_l2 & loguniform [1e-3, 100] & 65.3 \\
min\_gain\_to\_split & uniform [0, 1.0] & 0.092 \\
n\_estimators & early stop, max 1000 & 423 \\
\midrule
\multicolumn{3}{l}{\textbf{XGBoost \texttt{rank:ndcg}}} \\
\midrule
Hyperparameter & Search space & Selected \\
\midrule
eta & loguniform [0.01, 0.10] & 0.0687 \\
max\_depth & int [4, 12] & 9 \\
min\_child\_weight & int [1, 50] & 18 \\
subsample & uniform [0.5, 1.0] & 0.975 \\
colsample\_bytree & uniform [0.5, 1.0] & 0.765 \\
reg\_alpha & loguniform [1e-3, 100] & 7.50 \\
reg\_lambda & loguniform [1e-3, 100] & 99.2 \\
gamma & uniform [0, 5] & 1.63 \\
n\_estimators & early stop, max 1000 & 694 \\
\bottomrule
\end{tabular}
\end{table}

The selected configurations (Table~\ref{tab:hyper}) are informative on their own. The LightGBM best uses a large \texttt{num\_leaves} (110) combined with aggressive L2 regularisation (\texttt{lambda\_l2} = 65.3) and a tight \texttt{feature\_fraction} of 0.56. This is the classical bias-variance trade Optuna landed on: a high-capacity tree shape combined with strong regularisation in two complementary axes (column subsampling and L2 penalty). The XGBoost best uses a similar structure with a deeper tree (depth 9) and an even stronger L2 (\texttt{reg\_lambda} = 99.2) but no column subsampling pressure (\texttt{colsample\_bytree} = 0.77, looser). Both configurations therefore concentrate regularisation on the L2 axis rather than on tree depth or learning rate, which is consistent with the observation that the dataset has heavy categorical features (\texttt{prop\_id}, \texttt{srch\_destination\_id}) whose memorisation tendency benefits most from L2.

\begin{figure}[t]
\centering
\includegraphics[width=0.95\linewidth]{optuna_convergence.png}
\caption{Optuna TPE convergence over 50 trials. Each dot is one trial's val NDCG@5; the red curve is the best score so far. The LightGBM study (left) reaches its best at trial 38; the XGBoost study (right) converges more slowly with several near-best plateaus, suggesting a flatter optimum surface.}
\label{fig:optuna}
\end{figure}

Figure~\ref{fig:optuna} shows the convergence trajectories. The LightGBM study (left) stabilises by trial 30 and the warmup random trials are visibly below the TPE-driven trials, indicating that the sampler is correctly exploiting the surrogate model. The XGBoost study (right) finds its best near trial 22 and then plateaus, a sign that the L2-heavy region is a wide ridge rather than a sharp peak. The XGBoost best objective value from the trial (NDCG@5 = 0.40856 at 310 trees) is lower than the refit value of 0.41223 at 694 trees, because the refit uses early\_stopping = 50 rather than 30 and XGBoost was still improving slowly at trial 310. The downstream pipeline uses the refit value.

\subsection{Position-bias correction experiment}

About 30\% of training rows have \texttt{random\_bool} = 1, meaning Expedia served the search results in random order rather than ranked order. The standard fix for position bias in learning-to-rank is to upweight the rows where the position signal is decoupled from quality. We tested doubling the \texttt{sample\_weight} on \texttt{random\_bool} = 1 rows during training. The change hurt val NDCG@5 by 0.0042 on the Optuna-tuned LightGBM configuration. The plausible cause is that doubling the weight of 30\% of rows makes the model treat shuffled-page rows as more informative than ranked-page rows, but at inference time the test distribution is dominated by ranked pages. We reject the sample-weight approach. The fairness rerank in Section~\ref{sec:bias} attacks the same exposure-bias concern differently, as a deployment-time additive boost rather than a training-time loss reweighting, which avoids this trade-off.

\subsection{Seeded ensemble and LightGBM-XGBoost blend}\label{subsec:ensemble}

A three-seed score-averaged LightGBM ensemble (seeds 42, 123, 456) reaches val NDCG@5 = 0.40933, a +0.0028 lift over the best single seed (0.40999 at seed 42 minus 0.40629 at seed 123, balanced by 0.40651 at seed 456). The pairwise Spearman rank correlation between the three seed predictions on val sits near 0.99, which explains why the ensemble lift is modest: the seeds disagree on a small share of pairs. XGBoost on the same features and the same val slice reaches NDCG@5 = 0.41223 in the refit, which is +0.0029 over the LightGBM 3-seed ensemble and reflects a different point in the bias-variance trade.

The disagreement between LightGBM and XGBoost is the larger source of ensemble lift. On val the Spearman rank correlation between the LightGBM 3-seed average and XGBoost sits near 0.92, leaving enough disagreement that averaging the two implementations cancels correlated errors. We sweep the blend weight $w_{\text{xgb}} \in \{0, 0.25, 0.5, 0.6, 0.7, 0.75, 0.85, 1\}$ on val and select $w_{\text{xgb}}$ at the maximum, which lands at 0.5 with val NDCG@5 = 0.41498. The weight selection sits on a flat plateau (the 0.5 to 0.7 range differs by less than 0.0005 on val), which means the chosen weight is stable across the seed split rather than a sharp val-fit artefact.
